\documentclass{article}
\usepackage[T2A,russian]{graphicx, fontenc, amsmath,babel} % Required for inserting images
\title{182 уравнение гиперболического типа гл.11}

\begin{document}

\maketitle

плотностью тепловых источников F(x,t)F(x,t) в точке x в момент $t^1$.

В результате действия этих источников на участке стержня

(x,x+dx) за промежуток времени (t,t+dt) выделится количество тепла 
$$dQ = SF(x,t)dxdt\eqno(8)$$ 
или в интегральной форме
$$Q = S\int_{t_1}^{t_2}\int_{x_1}^{x_2}F(x,t)dxdt\eqno(9)$$
где Q — количество тепла, выделяющегося на участке стержня
(x1,x2) за промежуток времени (t1,t2).

Уравнение теплопроводности получается при подсчёте баланса тепла на некотором отрезке (x1,x2) за некоторый промежуток времени (t1,t2). Применяя закон сохранения энергии и пользуясь формулами (5), (7) и (9), можно написать равенство

\begin{multline}\tag{10}
\int_{t_1}^{t_2}[k\frac{\partial u}{\partial x}(x,\tau)|_{x=x_2}-k\frac{\partial u}{\partial x}(x,\tau)|_{x=x_1}]d\tau+\int_{x_1}^{x_3}\int_{t_1}^{t_\tau}F(\xi,\tau)d\xi d\tau= \\=\int_{x_1}^{x_3}cp[u(\xi,t_2)-u(\xi,t_1)]d\xi,
\end{multline}

которое и представляет уравнение теплопроводности в интегральной форме.

Чтобы получить уравнение теплопроводности в дифференциальной форме, предположим, что функция u(x,t) имеет непрерывные производные $u_{xx}$​ и $u_t$2)​.

Пользуясь теоремой о среднем, получаем равенство
\begin{multline}\tag{11}
[k\frac{\partial u}{\partial x}(x,\tau)|_{x=x_2}-k\frac{\partial u}{\partial x}(x,\tau)|_{x=x_1}]_{\tau=t_3}\Delta t+ F(x_4,t_4)\Delta x\Delta t\\=\{cp[u(\xi,t_2)-u(\xi,t_1)]\}_{\xi=x_3}\Delta x,
\end{multline}
\newpage
\begin{center}
\title{простейшие задачи            183}
\end{center}
которое при помощи теоремы о конечных приращениях можно преобразовать к виду
$$\frac{\partial}{\partial x}[k\frac{\partial u}{\partial x}(x,t)]_{x=x_5,t=t_3}\Delta t\Delta x + F(x_4,t_4)\Delta x\Delta t=[cp\frac{\partial u}{\partial t}(x,t)]_{x=x_3,t=t_5}\Delta x \Delta t, \eqno(12)$$
где $t_3,t_4,t_5$ и $x_3,x_4,x_5$ — промежуточные точки интервалов (t1,t2) и (x1,x2).

Отсюда, после сокращения на произведение ΔxΔtΔxΔt, находим:
$$\frac{\partial}{\partial x}(k\frac{\partial u}{\partial x})|_{x=x_5,t=t_3}+F(x,t)|_{x=x_4,t=t_4}=cp\frac{\partial u}{\partial t}|_{t=t_5,x=x_3}\eqno(13)$$
Все эти рассуждения относятся к произвольным промежуткам (x1,x2) и (t1,t2). Переходя к пределу при $x1,x2→x$ и $t1,t2→t$, получим уравнение
$$\frac{\partial}{\partial x}(k\frac{\partial u}{\partial x})+F(x,t)=cp\frac{\partial u}{\partial t},\eqno(14)$$
называемое уравнением теплопроводности.

Рассмотрим некоторые частные случаи.

1.  Если стержень однороден, то k,c,p можно считать постоянными, и уравнение обычно записывают ввиде
$$u_t = a^2u_{xx}+f(x,t),$$
$$a^2=\frac{k}{cp},\ \ f(x,t)=\frac{F(x,t}{cp},$$
где $a^2$ — постоянная, называется коэффициентом температуропроводности. Если источники отсутствуют, т. е. F(x,t)=0, то уравнение теплопроводности принимает простой вид:
$$u_t=a^2u_{xx}.\eqno(14')$$
Плотность тепловых источников может зависеть от температуры. В случае теплообмена с окружающей средой, подчиняющегося закону Ньютона, количество тепла, теряемого стержнем1), рассчитываемое на единицу длины и времени, равно

$$F_0=h(u-\theta),$$

\end{document}
